% GENERATED FILE. Edit canonical lessons, then run npm run generate:books.
% canonical-source-hash: fnv1a64:de9347dd2b5575e1
% canonical-lessons: GE-C312-volksfest, GE-C312-jahrmarkt, GE-C312-karussell, GE-C312-tracht, GE-C312-kapelle

\chapter{Das Volksfest}
\label{ch:ge-das-volksfest}
\hlchaptermodality{312}

\begin{chapteropening}
\textbf{By the end of this chapter:} \emph{I can talk about a fair: a public festival, a funfair, a merry-go-round, traditional dress and a band.}

Complete the last lesson of chapter 312: I can talk about a fair: a public festival, a funfair, a merry-go-round, traditional dress and a band.
\end{chapteropening}

\section[das Volksfest]{das Volksfest — the public festival, the folk festival}

\label{lesson:GE-C312-volksfest}



\begin{quote}
Before the new one: say the German for the giving of Christmas presents, then the German for the gingerbread.
\end{quote}

\subsection*{You'll want to know: das Volksfest}
\textbf{das Volksfest} — \textquotedblleft{}the public festival, the folk festival\textquotedblright{}.

\begin{quote}
Im Herbst gibt es in unserer Stadt ein großes Volksfest.
\end{quote}

\begin{quote}
\emph{In autumn there is a big public festival in our town.}
\end{quote}

Plural: \emph{die Volksfeste}.

\begin{cousinweb}
\emph{Volksfest} joins \emph{Volk} \textquotedblleft{}people\textquotedblright{}, Old High German \emph{folc} (EtymDB), and \emph{Fest} \textquotedblleft{}festival\textquotedblright{}, from Latin \emph{festum} \textquotedblleft{}feast\textquotedblright{}.
\end{cousinweb}

\begin{grammarlens}[title={What you've built}]
One more word for telling, explaining and describing at B1.
\end{grammarlens}

\subsection*{Your turn}
\begin{itemize}
\raggedright
  \item \emph{Say it:} \emph{das Volksfest}
  \item \emph{Say it:} \emph{das Volksfest}, once more
  \item \emph{Say it:} a sentence of your own with \emph{das Volksfest}
  \item \emph{Recall:} say the German for the giving of Christmas presents, then the German for the gingerbread, then say \emph{das Volksfest} again
\end{itemize}

\begin{tcolorbox}[breakable,colback=teal!4,colframe=teal!35!black,title={Before you move on}]
What does das Volksfest mean? (\textquotedblleft{}The public festival, the folk festival\textquotedblright{}.) Say it once more.
\end{tcolorbox}
\section[der Jahrmarkt]{der Jahrmarkt — the funfair, the fair}

\label{lesson:GE-C312-jahrmarkt}



\begin{quote}
Before the new one: say the German for the gingerbread, then the German for the public festival.
\end{quote}

\subsection*{You'll want to know: der Jahrmarkt}
\textbf{der Jahrmarkt} — \textquotedblleft{}the funfair, the fair\textquotedblright{}.

\begin{quote}
Auf dem Jahrmarkt aßen wir gebrannte Mandeln.
\end{quote}

\begin{quote}
\emph{At the fair we ate roasted almonds.}
\end{quote}

Plural: \emph{die Jahrmärkte}.

\begin{cousinweb}
\emph{Jahrmarkt} joins \emph{Jahr} \textquotedblleft{}year\textquotedblright{} and \emph{Markt} \textquotedblleft{}market\textquotedblright{}, from Latin \emph{mercātus} (EtymDB): a market held once a year.
\end{cousinweb}

\begin{grammarlens}[title={What you've built}]
One more word for telling, explaining and describing at B1.
\end{grammarlens}

\subsection*{Your turn}
\begin{itemize}
\raggedright
  \item \emph{Say it:} \emph{der Jahrmarkt}
  \item \emph{Say it:} \emph{der Jahrmarkt}, once more
  \item \emph{Say it:} a sentence of your own with \emph{der Jahrmarkt}
  \item \emph{Recall:} say the German for the gingerbread, then the German for the public festival, then say \emph{der Jahrmarkt} again
\end{itemize}

\begin{tcolorbox}[breakable,colback=teal!4,colframe=teal!35!black,title={Before you move on}]
What does der Jahrmarkt mean? (\textquotedblleft{}The funfair, the fair\textquotedblright{}.) Say it once more.
\end{tcolorbox}
\section[das Karussell]{das Karussell — the merry-go-round, the carousel}

\label{lesson:GE-C312-karussell}



\begin{quote}
Before the new one: say the German for the public festival, then the German for the funfair.
\end{quote}

\subsection*{You'll want to know: das Karussell}
\textbf{das Karussell} — \textquotedblleft{}the merry-go-round, the carousel\textquotedblright{}.

\begin{quote}
Die Kinder fuhren dreimal Karussell.
\end{quote}

\begin{quote}
\emph{The children went on the merry-go-round three times.}
\end{quote}

Plural: \emph{die Karussells}. \emph{Karussell fahren}.

\begin{cousinweb}
\emph{Karussell} was borrowed from French \emph{carrousel}, which came from Italian \emph{carosello} (EtymDB).
\end{cousinweb}

\begin{grammarlens}[title={What you've built}]
One more word for telling, explaining and describing at B1.
\end{grammarlens}

\subsection*{Your turn}
\begin{itemize}
\raggedright
  \item \emph{Say it:} \emph{das Karussell}
  \item \emph{Say it:} \emph{das Karussell}, once more
  \item \emph{Say it:} a sentence of your own with \emph{das Karussell}
  \item \emph{Recall:} say the German for the public festival, then the German for the funfair, then say \emph{das Karussell} again
\end{itemize}

\begin{tcolorbox}[breakable,colback=teal!4,colframe=teal!35!black,title={Before you move on}]
What does das Karussell mean? (\textquotedblleft{}The merry-go-round, the carousel\textquotedblright{}.) Say it once more.
\end{tcolorbox}
\section[die Tracht]{die Tracht — the traditional dress}

\label{lesson:GE-C312-tracht}



\begin{quote}
Before the new one: say the German for the funfair, then the German for the merry-go-round.
\end{quote}

\subsection*{You'll want to know: die Tracht}
\textbf{die Tracht} — \textquotedblleft{}the traditional dress\textquotedblright{}.

\begin{quote}
Viele Leute trugen Tracht.
\end{quote}

\begin{quote}
\emph{Many people wore traditional dress.}
\end{quote}

Plural: \emph{die Trachten}. In the south and in Austria, the \emph{Dirndl} and the \emph{Lederhose}.

\begin{cousinweb}
\emph{Tracht} is inherited from Old High German \emph{draht}, built on the root of \emph{tragen} \textquotedblleft{}to carry, to wear\textquotedblright{}, Proto-Germanic \emph{draganą} (EtymDB): what one wears.
\end{cousinweb}

\begin{grammarlens}[title={What you've built}]
One more word for telling, explaining and describing at B1.
\end{grammarlens}

\subsection*{Your turn}
\begin{itemize}
\raggedright
  \item \emph{Say it:} \emph{die Tracht}
  \item \emph{Say it:} \emph{die Tracht}, once more
  \item \emph{Say it:} a sentence of your own with \emph{die Tracht}
  \item \emph{Recall:} say the German for the funfair, then the German for the merry-go-round, then say \emph{die Tracht} again
\end{itemize}

\begin{tcolorbox}[breakable,colback=teal!4,colframe=teal!35!black,title={Before you move on}]
What does die Tracht mean? (\textquotedblleft{}The traditional dress\textquotedblright{}.) Say it once more.
\end{tcolorbox}
\section[die Kapelle]{die Kapelle — the band (of musicians); the chapel}

\label{lesson:GE-C312-kapelle}



\begin{quote}
Before the new one: say the German for the merry-go-round, then the German for the traditional dress.
\end{quote}

\subsection*{You'll want to know: die Kapelle}
\textbf{die Kapelle} — \textquotedblleft{}the band (of musicians); the chapel\textquotedblright{}.

\begin{quote}
Eine Kapelle spielte alte Lieder.
\end{quote}

\begin{quote}
\emph{A band played old songs.}
\end{quote}

Plural: \emph{die Kapellen}. Also a small church, a chapel.

\begin{cousinweb}
\emph{Kapelle} comes from Middle High German \emph{kappelle} (EtymDB), from Medieval Latin \emph{capella} \textquotedblleft{}a little cloak\textquotedblright{}, then, as the story is usually told, the chapel where a famous cloak was kept, then its musicians.
\end{cousinweb}

\begin{grammarlens}[title={What you've built}]
One more word for telling, explaining and describing at B1.
\end{grammarlens}

\subsection*{Your turn}
\begin{itemize}
\raggedright
  \item \emph{Say it:} \emph{die Kapelle}
  \item \emph{Say it:} \emph{die Kapelle}, once more
  \item \emph{Say it:} a sentence of your own with \emph{die Kapelle}
  \item \emph{Recall:} say the German for the merry-go-round, then the German for the traditional dress, then say \emph{die Kapelle} again
\end{itemize}

\begin{tcolorbox}[breakable,colback=teal!4,colframe=teal!35!black,title={Before you move on}]
What does die Kapelle mean? (\textquotedblleft{}The band (of musicians); the chapel\textquotedblright{}.) Say it once more.
\end{tcolorbox}
